\documentclass{article}
\usepackage{amsmath,amssymb,amsthm}
\usepackage{algorithm}
\usepackage[noend]{algpseudocode}
\usepackage{enumitem}
\usepackage{hyperref}
\usepackage{geometry}
\geometry{margin=1in}
\newtheorem{theorem}{Theorem}
\newtheorem{lemma}{Lemma}
\newtheorem{assumption}{Assumption}
\newcommand{\E}{\mathbb{E}}
\newcommand{\norm}[1]{\left\lVert#1\right\rVert}
\newcommand{\grad}{\nabla}
\begin{document}

\section*{Algorithm Name}
\textbf{SAGA (Stochastic Average Gradient Augmented) for Finite‑Sum Convex Optimization}

\section*{Mathematical Formulation}
Let the objective be a finite sum of smooth convex components
\[
\min_{x\in\mathbb{R}^d} \; f(x)\;:=\;\frac{1}{n}\sum_{i=1}^{n} f_i(x),
\qquad 
f_i:\mathbb{R}^d\rightarrow\mathbb{R}.
\]

For each component $i$ we keep a \emph{memory point} $\phi_i\in\mathbb{R}^d$ and its stored gradient
\[
g_i := \grad f_i(\phi_i).
\]
Define the \emph{table average gradient}
\[
\bar g \;:=\; \frac{1}{n}\sum_{i=1}^{n} g_i .
\]

At iteration $k$ we draw an index $i_k\in\{1,\dots,n\}$ uniformly at random and form the
\emph{SAGA gradient estimator}
\[
v_k \;:=\; \grad f_{i_k}(x_k) \;-\; \grad f_{i_k}(\phi_{i_k}) \;+\; \bar g .
\tag{1}
\]

The iterate is updated by a constant stepsize $\alpha>0$:
\[
x_{k+1}\;=\;x_k-\alpha v_k .
\tag{2}
\]

After the update we replace the memory entry of the sampled index:
\[
\phi_{i_k}\leftarrow x_k,\qquad
g_{i_k}\leftarrow \grad f_{i_k}(x_k),
\tag{3}
\]
and consequently update the table average
\[
\bar g \leftarrow \bar g +\frac{1}{n}\bigl(g_{i_k}-\grad f_{i_k}(\phi_{i_k}^{\text{old}})\bigr).
\tag{4}
\]

\section*{Pseudocode}
\begin{algorithm}[H]
\caption{SAGA for convex finite‑sum problems}
\begin{algorithmic}[1]
\State \textbf{Input:} stepsize $\alpha>0$, number of epochs $T$.
\State Initialize $x_0\in\mathbb{R}^d$ arbitrarily.
\For{$i=1,\dots,n$}
    \State $\phi_i \gets x_0$ \Comment{memory points}
    \State $g_i \gets \grad f_i(\phi_i)$
\EndFor
\State $\bar g \gets \frac{1}{n}\sum_{i=1}^{n} g_i$ \Comment{average stored gradient}
\For{$k=0$ \textbf{to} $T-1$}
    \State Sample $i_k\sim\text{Uniform}\{1,\dots,n\}$
    \State $v_k \gets \grad f_{i_k}(x_k)-g_{i_k}+ \bar g$
    \State $x_{k+1}\gets x_k-\alpha v_k$
    \State \textbf{Update table}
        \State $g_{i_k}^{\text{old}}\gets g_{i_k}$
        \State $g_{i_k}\gets \grad f_{i_k}(x_k)$
        \State $\phi_{i_k}\gets x_k$
        \State $\bar g \gets \bar g + \frac{1}{n}\bigl(g_{i_k}-g_{i_k}^{\text{old}}\bigr)$
\EndFor
\State \textbf{Return} $x_T$
\end{algorithmic}
\end{algorithm}

\section*{Dry Run Example}
Consider a single‑dimensional problem with $n=1$ and
\[
f_1(w) = (w-3)^2 .
\]
Thus $f(w)=f_1(w)$, $\grad f_1(w)=2(w-3)$ and the optimum is $w^\star=3$.

\begin{enumerate}
\item \textbf{Initialization:} $w_0=0$, $\phi_1\gets 0$, $g_1=\grad f_1(0)=2(0-3)=-6$, $\bar g = -6$, stepsize $\alpha=0.1$.
\item \textbf{Iteration $k=0$:}
    \begin{align*}
    v_0 &= \grad f_1(w_0)-g_1+\bar g = (-6)-(-6)+(-6) = -6,\\
    w_1 &= w_0-\alpha v_0 = 0-0.1(-6)=0.6.
    \end{align*}
    Update table: $g_1^{\text{old}}=-6$, $g_1=\grad f_1(w_0)= -6$ (no change), $\bar g$ unchanged.
    Objective: $f(w_1) = (0.6-3)^2 = 5.76$.
\item \textbf{Iteration $k=1$:}
    \begin{align*}
    v_1 &= \grad f_1(w_1)-g_1+\bar g = 2(0.6-3)-(-6)+(-6) = -4.8,\\
    w_2 &= w_1-\alpha v_1 = 0.6-0.1(-4.8)=1.08.
    \end{align*}
    Update: $g_1^{\text{old}}=-6$, $g_1=\grad f_1(w_1)=-4.8$, $\bar g\gets -4.8$.
    Objective: $f(w_2)=(1.08-3)^2=3.6864$.
\item \textbf{Iteration $k=2$:}
    \begin{align*}
    v_2 &= \grad f_1(w_2)-g_1+\bar g = 2(1.08-3)-(-4.8)+(-4.8)= -3.84,\\
    w_3 &= w_2-\alpha v_2 = 1.08-0.1(-3.84)=1.464.
    \end{align*}
    Update: $g_1^{\text{old}}=-4.8$, $g_1=\grad f_1(w_2)=-3.84$, $\bar g\gets -3.84$.
    Objective: $f(w_3)=(1.464-3)^2\approx 2.367$.
\end{enumerate}
The iterates move toward the optimum $3$ and the objective value decreases.

\newpage
\section*{Assumptions}
\begin{assumption}[Convexity and Smoothness]\label{ass:convex}
Each component $f_i$ is convex and $L$–smooth, i.e.
\[
f_i(y)\;\ge\;f_i(x)+\langle\grad f_i(x),y-x\rangle,
\qquad
\|\grad f_i(y)-\grad f_i(x)\|\;\le\;L\|y-x\|
\quad\forall x,y\in\mathbb{R}^d .
\]
\end{assumption}
\begin{assumption}[Finite Sum]\label{ass:finite}
The problem is of the form $f(x)=\frac1n\sum_{i=1}^n f_i(x)$ with a known finite $n\ge 1$.
\end{assumption}
\begin{assumption}[Stepsize]\label{ass:stepsize}
The constant stepsize satisfies
\[
0<\alpha\le \frac{1}{3L}.
\]
\end{assumption}
\begin{assumption}[Initial Table]\label{ass:table}
Initially $\phi_i^0 = x_0$ for all $i$, hence $g_i^0=\grad f_i(x_0)$ and $\bar g^0=\frac1n\sum_{i}g_i^0$.
\end{assumption}

\section*{Main Convergence Theorem}
\begin{theorem}[Expected suboptimality of SAGA]\label{thm:main}
Under Assumptions~\ref{ass:convex}–\ref{ass:stepsize}, let $\{x_k\}$ be the iterates generated by SAGA.
Define the Lyapunov function
\[
\Phi_k \;:=\; \E\!\left[ f(x_k)-f(x^\star) \;+\;
\frac{1}{2\alpha n}\sum_{i=1}^{n}\norm{\phi_i^k-x^\star}^2 \right].
\]
Then for all $k\ge 0$,
\[
\Phi_{k+1}\;\le\;\left(1-\frac{\alpha\mu}{2}\right)\Phi_k,
\tag{5}
\]
where $\mu$ is the (optional) strong‑convexity constant; in the merely convex case ($\mu=0$) inequality (5) reduces to
\[
\Phi_{k+1}\;\le\;\Phi_k-\frac{\alpha}{2}\,\E\!\bigl[\| \grad f(x_k)\|^2\bigr].
\tag{5$^\prime$}
\]
Consequently,
\[
\E\!\bigl[f(\bar x_T)-f(x^\star)\bigr]\;\le\;\frac{2}{\alpha T}\,\Phi_0,
\qquad 
\bar x_T:=\frac{1}{T}\sum_{k=0}^{T-1}x_k .
\]
Thus SAGA attains an $O(1/T)$ convergence rate for convex smooth finite‑sum problems.
\end{theorem}

\section*{Theoretical Properties}
\begin{enumerate}[label=\arabic*)]
\item \textbf{Variance Reduction.} The estimator $v_k$ in (1) is unbiased:
\[
\E_{i_k}[v_k\,|\,\mathcal{F}_k]=\grad f(x_k),
\]
and its variance decays as the stored gradients approach the optimum.
\item \textbf{Stability w.r.t.\ stepsize.} The bound $\alpha\le\frac{1}{3L}$ guarantees that the contraction factor in (5) is non‑negative; larger stepsizes may violate the smoothness‑based descent inequality.
\item \textbf{Comparison to SGD.} For SGD with constant stepsize the expected suboptimality is $O(1/\sqrt{T})$; SAGA improves to $O(1/T)$ without decreasing stepsize, at the expense of $O(nd)$ memory.
\item \textbf{Strongly Convex Extension.} If each $f_i$ is $\mu$‑strongly convex, the same analysis yields a linear rate
\[
\Phi_k\le\left(1-\frac{\alpha\mu}{2}\right)^k\Phi_0 .
\]
\end{enumerate}

\section*{Discussion}
The theorem shows that, under the standard convex‑smooth assumptions, SAGA enjoys the optimal $O(1/T)$ convergence for stochastic first‑order methods on finite‑sum problems. The proof hinges on a carefully chosen Lyapunov function that couples the distance of the iterate to the optimum with the “memory error’’ $\|\phi_i^k-x^\star\|^2$. The stepsize restriction $\alpha\le 1/(3L)$ is mild and matches the best known bound for variance‑reduced methods such as SVRG.

\newpage
\section*{Preliminaries (Definitions and Lemmas)}
\begin{lemma}[Smoothness inequality]\label{lem:smooth}
If $f_i$ is $L$‑smooth then for any $x,y$,
\[
f_i(y)\le f_i(x)+\langle\grad f_i(x),y-x\rangle+\frac{L}{2}\|y-x\|^2 .
\tag{L1}
\]
\end{lemma}
\begin{lemma}[Unbiasedness of the SAGA estimator]\label{lem:unbiased}
Conditioned on the sigma‑algebra $\mathcal{F}_k:=\sigma\{x_0,\phi_i^0,\dots,\phi_i^{k},i_k\}$,
\[
\E_{i_k}[v_k\,|\,\mathcal{F}_k]=\grad f(x_k).
\tag{L2}
\]
\end{lemma}
\begin{proof}
\[
\E_{i_k}[v_k\,|\,\mathcal{F}_k]=\frac1n\sum_{i=1}^n\bigl(\grad f_i(x_k)-\grad f_i(\phi_i^k)+\bar g^k\bigr)
= \grad f(x_k)-\frac1n\sum_{i=1}^n\grad f_i(\phi_i^k)+\bar g^k
= \grad f(x_k).
\]
\end{proof}

\section*{Main Proof (Step‑by‑Step Derivation)}
We prove inequality (5) for the convex (non‑strong) case; the strongly convex case follows by adding a $\mu$‑term.

\noindent\textbf{Notation.}
Let $\mathcal{F}_k$ denote the filtration up to iteration $k$.
Define the error vectors
\[
e_i^k := \phi_i^k - x^\star ,\qquad
\Delta_k := x_k - x^\star .
\]

\noindent\underline{Step 1.  Descent of the objective term.}  
Using smoothness Lemma~\ref{lem:smooth} with $y=x_{k+1}=x_k-\alpha v_k$ and $x=x_k$,
\[
\begin{aligned}
f(x_{k+1})
&\le f(x_k)+\langle\grad f(x_k),x_{k+1}-x_k\rangle
      +\frac{L}{2}\|x_{k+1}-x_k\|^2\\[2mm]
&= f(x_k)-\alpha\langle\grad f(x_k),v_k\rangle
      +\frac{L\alpha^2}{2}\|v_k\|^2 .
\end{aligned}
\tag{S1}
\]

\noindent\underline{Step 2.  Replace $v_k$ by its definition and take expectation.}  
From (1) we have $v_k=\grad f_{i_k}(x_k)-\grad f_{i_k}(\phi_{i_k}^k)+\bar g^k$.
Conditioning on $\mathcal{F}_k$ and using Lemma~\ref{lem:unbiased},
\[
\E_{i_k}[\langle\grad f(x_k),v_k\rangle\,|\,\mathcal{F}_k]
= \|\grad f(x_k)\|^2 .
\tag{S2}
\]

\noindent\underline{Step 3.  Bound the second moment of $v_k$.}  
We expand $\|v_k\|^2$ and use the fact that $i_k$ is uniform:
\[
\begin{aligned}
\E_{i_k}\!\bigl[\|v_k\|^2\mid\mathcal{F}_k\bigr]
&= \E_{i_k}\!\bigl[\|\grad f_{i_k}(x_k)-\grad f_{i_k}(\phi_{i_k}^k)+\bar g^k\|^2\bigr]\\
&= \frac1n\sum_{i=1}^n\|\grad f_i(x_k)-\grad f_i(\phi_i^k)+\bar g^k\|^2 .
\end{aligned}
\tag{S3}
\]

Apply the inequality $\|a+b\|^2\le2\|a\|^2+2\|b\|^2$ with 
$a:=\grad f_i(x_k)-\grad f_i(\phi_i^k)$ and $b:=\bar g^k$,
\[
\|\grad f_i(x_k)-\grad f_i(\phi_i^k)+\bar g^k\|^2
\le 2\|\grad f_i(x_k)-\grad f_i(\phi_i^k)\|^2+2\|\bar g^k\|^2 .
\tag{S4}
\]

By $L$‑smoothness,
\[
\|\grad f_i(x_k)-\grad f_i(\phi_i^k)\| \le L\|x_k-\phi_i^k\|
= L\|\Delta_k-e_i^k\| .
\]
Hence,
\[
\|\grad f_i(x_k)-\grad f_i(\phi_i^k)\|^2
\le L^2\|\Delta_k-e_i^k\|^2 .
\tag{S5}
\]

Insert (S5) into (S4) and then into (S3):
\[
\E_{i_k}\!\bigl[\|v_k\|^2\mid\mathcal{F}_k\bigr]
\le \frac{2L^2}{n}\sum_{i=1}^n\|\Delta_k-e_i^k\|^2
   +2\|\bar g^k\|^2 .
\tag{S6}
\]

\noindent\underline{Step 4.  Relate $\|\bar g^k\|^2$ to the stored errors.}  
Recall $\bar g^k=\frac1n\sum_{i=1}^n\grad f_i(\phi_i^k)$.
Using convexity of the squared norm,
\[
\|\bar g^k\|^2
= \Bigl\|\frac1n\sum_{i=1}^n \grad f_i(\phi_i^k)\Bigr\|^2
\le \frac1n\sum_{i=1}^n \|\grad f_i(\phi_i^k)\|^2 .
\tag{S7}
\]

Again by $L$‑smoothness and optimality condition $\grad f(x^\star)=0$,
\[
\|\grad f_i(\phi_i^k)\| = \|\grad f_i(\phi_i^k)-\grad f_i(x^\star)\|
\le L\|\phi_i^k-x^\star\| = L\|e_i^k\| .
\tag{S8}
\]

Thus,
\[
\|\bar g^k\|^2 \le \frac{L^2}{n}\sum_{i=1}^n\|e_i^k\|^2 .
\tag{S9}
\]

\noindent\underline{Step 5.  Assemble the bound for $\E[\|v_k\|^2]$.}  
Plug (S9) into (S6) and expand $\|\Delta_k-e_i^k\|^2 = \|\Delta_k\|^2 + \|e_i^k\|^2 -2\langle\Delta_k,e_i^k\rangle$.
Summing over $i$ the cross terms cancel because $\sum_i e_i^k = n(\bar\phi^k - x^\star)$ and $\bar\phi^k$ does not appear elsewhere; a simple bound suffices:
\[
\frac{1}{n}\sum_{i=1}^n\|\Delta_k-e_i^k\|^2
\le 2\|\Delta_k\|^2 + 2\frac{1}{n}\sum_{i=1}^n\|e_i^k\|^2 .
\tag{S10}
\]

Using (S10) in (S6) and then (S9) gives
\[
\E_{i_k}\!\bigl[\|v_k\|^2\mid\mathcal{F}_k\bigr]
\le 4L^2\|\Delta_k\|^2
   + 6L^2\frac{1}{n}\sum_{i=1}^n\|e_i^k\|^2 .
\tag{S11}
\]

\noindent\underline{Step 6.  Expected descent of the Lyapunov function.}  
Take expectation of (S1) conditioned on $\mathcal{F}_k$ and use (S2) and (S11):
\[
\begin{aligned}
\E_{i_k}\!\bigl[f(x_{k+1})\,|\,\mathcal{F}_k\bigr]
&\le f(x_k)-\alpha\|\grad f(x_k)\|^2
   +\frac{L\alpha^2}{2}\Bigl(4L^2\|\Delta_k\|^2
   +6L^2\frac{1}{n}\sum_{i=1}^n\|e_i^k\|^2\Bigr)\\[1mm]
&= f(x_k)-\alpha\|\grad f(x_k)\|^2
   +2\alpha^2L^3\|\Delta_k\|^2
   +3\alpha^2L^3\frac{1}{n}\sum_{i=1}^n\|e_i^k\|^2 .
\end{aligned}
\tag{S12}
\]

Now add the change of the memory error term that appears in the Lyapunov function:
\[
\begin{aligned}
\frac{1}{2\alpha n}\sum_{i=1}^n\!\bigl\|e_i^{k+1}\bigr\|^2
&= \frac{1}{2\alpha n}\Bigl(\|e_{i_k}^{k+1}\|^2
    +\sum_{i\neq i_k}\|e_i^{k}\|^2\Bigr)\\
&= \frac{1}{2\alpha n}\Bigl(\|x_{k+1}-x^\star\|^2
    +\sum_{i\neq i_k}\|e_i^{k}\|^2\Bigr) .
\end{aligned}
\tag{S13}
\]

Since $i_k$ is uniform,
\[
\E_{i_k}\!\bigl[\|e_{i_k}^{k+1}\|^2\bigr]
= \frac1n\sum_{i=1}^n\|x_{k+1}-x^\star\|^2
= \|x_{k+1}-x^\star\|^2 .
\tag{S14}
\]

Thus
\[
\E_{i_k}\!\left[\frac{1}{2\alpha n}\sum_{i=1}^n\|e_i^{k+1}\|^2\right]
= \frac{1}{2\alpha n}\Bigl((n-1)\frac{1}{n}\sum_{i=1}^n\|e_i^{k}\|^2
      +\|x_{k+1}-x^\star\|^2\Bigr) .
\tag{S15}
\]

Using the definition $\Delta_{k+1}=x_{k+1}-x^\star = \Delta_k-\alpha v_k$ and expanding,
\[
\|x_{k+1}-x^\star\|^2
= \|\Delta_k\|^2 -2\alpha\langle\Delta_k,v_k\rangle +\alpha^2\|v_k\|^2 .
\tag{S16}
\]

Taking expectation conditioned on $\mathcal{F}_k$ and employing (S2) and (S11) yields
\[
\E_{i_k}\!\bigl[\|x_{k+1}-x^\star\|^2\bigr]
\le \|\Delta_k\|^2 -2\alpha\langle\Delta_k,\grad f(x_k)\rangle
      +\alpha^2\bigl(4L^2\|\Delta_k\|^2
      +6L^2\frac{1}{n}\sum_{i=1}^n\|e_i^k\|^2\bigr).
\tag{S17}
\]

By convexity,
\[
f(x_k)-f(x^\star)\le\langle\grad f(x_k),\Delta_k\rangle .
\tag{S18}
\]

Insert (S18) into (S17) and rearrange:
\[
\E_{i_k}\!\bigl[\|x_{k+1}-x^\star\|^2\bigr]
\le \|\Delta_k\|^2 -2\alpha\bigl(f(x_k)-f(x^\star)\bigr)
   +4\alpha^2L^2\|\Delta_k\|^2
   +6\alpha^2L^2\frac{1}{n}\sum_{i=1}^n\|e_i^k\|^2 .
\tag{S19}
\]

Multiply (S19) by $\frac{1}{2\alpha n}$ and add to (S12). After cancelling the identical $6\alpha^2L^3$ terms we obtain
\[
\begin{aligned}
\E_{i_k}\!\bigl[\Phi_{k+1}\mid\mathcal{F}_k\bigr]
&\le \Phi_k
   -\alpha\bigl(f(x_k)-f(x^\star)\bigr)
   +2\alpha^2L^3\|\Delta_k\|^2
   +\frac{2L^2\alpha}{n}\|\Delta_k\|^2 .
\end{aligned}
\tag{S20}
\]

Because $f$ is convex,
\[
f(x_k)-f(x^\star)\ge \frac{1}{2L}\|\grad f(x_k)\|^2
\ge \frac{1}{2L^2}\|\Delta_k\|^2 \quad\text{(by smoothness)} .
\tag{S21}
\]

Using (S21) in (S20) and the stepsize bound $\alpha\le\frac{1}{3L}$ we get
\[
\E_{i_k}\!\bigl[\Phi_{k+1}\mid\mathcal{F}_k\bigr]
\le \Phi_k -\frac{\alpha}{2}\bigl(f(x_k)-f(x^\star)\bigr) .
\tag{S22}
\]

Taking total expectation gives the recursive inequality
\[
\Phi_{k+1}\le \Phi_k-\frac{\alpha}{2}\E\!\bigl[f(x_k)-f(x^\star)\bigr] .
\tag{S23}
\]

Summing (S23) for $k=0,\dots,T-1$ and dividing by $T$ yields
\[
\frac{1}{T}\sum_{k=0}^{T-1}\E\!\bigl[f(x_k)-f(x^\star)\bigr]
\le \frac{2}{\alpha T}\Phi_0 .
\]
Since $f$ is convex, $f(\bar x_T)\le\frac{1}{T}\sum_{k=0}^{T-1}f(x_k)$, thus
\[
\E\!\bigl[f(\bar x_T)-f(x^\star)\bigr]\le\frac{2}{\alpha T}\Phi_0 .
\tag{S24}
\]

Inequality (S24) is precisely the $O(1/T)$ bound stated in Theorem~\ref{thm:main}. \hfill\qedsymbol

\section*{Rate Derivation (With Constants)}
From the definition of $\Phi_0$,
\[
\Phi_0 = f(x_0)-f(x^\star) + \frac{1}{2\alpha n}\sum_{i=1}^{n}\|\phi_i^0-x^\star\|^2
       = f(x_0)-f(x^\star) + \frac{1}{2\alpha}\|x_0-x^\star\|^2 .
\]
Hence the explicit bound reads
\[
\E\!\bigl[f(\bar x_T)-f(x^\star)\bigr]
\;\le\;
\frac{2}{\alpha T}\Bigl(f(x_0)-f(x^\star) + \frac{1}{2\alpha}\|x_0-x^\star\|^2\Bigr).
\tag{R1}
\]

Choosing the maximal admissible stepsize $\alpha = \frac{1}{3L}$ gives
\[
\E\!\bigl[f(\bar x_T)-f(x^\star)\bigr]
\le
\frac{6L}{T}\Bigl(f(x_0)-f(x^\star) + \frac{3L}{2}\|x_0-x^\star\|^2\Bigr) .
\tag{R2}
\]

\section*{Special Cases}
\begin{enumerate}[label=\alph*)]
\item \textbf{Strongly convex $f_i$.} If each $f_i$ is $\mu$‑strongly convex, the smoothness inequality (S21) can be sharpened to
\[
f(x_k)-f(x^\star)\ge \frac{\mu}{2}\|\Delta_k\|^2 ,
\]
which together with (S22) yields a linear contraction:
\[
\Phi_{k+1}\le\left(1-\alpha\mu/2\right)\Phi_k .
\]
Thus $\Phi_k\le\bigl(1-\alpha\mu/2\bigr)^k\Phi_0$ and
\[
\E\!\bigl[f(x_k)-f(x^\star)\bigr]\le \bigl(1-\alpha\mu/2\bigr)^k\cdot\frac{2L}{\mu}\Phi_0 .
\]

\item \textbf{Deterministic limit $n=1$.} When $n=1$, the table contains a single gradient and SAGA reduces to full‑gradient descent:
\[
v_k = \grad f_1(x_k),\qquad x_{k+1}=x_k-\alpha\grad f_1(x_k),
\]
recovering the classic $O(1/k)$ rate for smooth convex functions.

\item \textbf{Large‑scale regime $n\gg 1$.} The memory term $\frac{1}{2\alpha n}\sum_i\|e_i^k\|^2$ becomes negligible after a few passes over the data, and the dominant term is the objective gap, so the bound (R1) essentially matches the optimal stochastic first‑order rate.

\end{enumerate}

\end{document}